\section{Conclusion}
\label{section:conclusion}
In this paper, In this paper, we present a novel iterative solver for the nonlinear system at each time step arising from the application of the convex splitting scheme to the Allen-Cahn equation with Flory–Huggins potential. 
We employ a ADMM framework to solve the resulting problem 
by reformulating the nonlinear system equivalently as a convex optimization problem for finding a minimizer. 
For theoretical part, on the one hand, we analyze the well-posedness and discrete energy stability of the convex splitting scheme; 
on the other hand, what is more critical is that we establish the unconditional convergence and bound-preserving property of the iterative algorithm through rigorous proof.
The effectiveness and robustness of the algorithm are demonstrated through three numerical examples, 
which serve to indirectly confirm the theoretical analysis.
\par The design philosophy for the iterative solver to the second-order convex splitting scheme follows essentially the same line of reasoning \cite{li2025overcoming}, 
yet one may fail to obtain discrete energy stability corresponding to the original energy \cite{gomez2011provably,yang2016linear}. 
What's more, while the current work focuses on the Allen–Cahn equation with periodic boundary conditions, 
the proposed solver framework can be readily extended to other boundary condition types, such as Neumann and Dirichlet types, 
and to more general spatial discretizations, such as the finite element method \cite{yuan2022second}.
\par Despite the promising results, several directions merit further investigation. 
(a) The adaptive adjustment strategy for the penalty parameter $\rho$
used in our experiments, although effective, could be further optimized.
(b) Constructing effective initial guesses for the ADMM iterations, 
possibly based on extrapolation techniques from previous time steps or on solutions from coarser grids, 
may lead to substantial reductions in iteration counts, especially for long-time simulations where the solution evolves slowly.
(c) Our solver is tailored to the convex splitting scheme. 
A natural direction for future research is to explore the applicability of this approach to other time-stepping methods.
